% Section 8.3: The Poisson Equation and Scale Invariance
% Complete LaTeX content following WHY → WHAT → SO WHAT framework

\subsection{The Poisson Equation and Scale Invariance}
\label{subsec:poisson-scale-invariance}

% ============================================================
% WHY: MOTIVATION (The Scaling Problem)
% ============================================================

We have established the gravitational Poisson equation at a single scale. But gravity must work at \emph{all} scales: from atoms to galaxies. A star contains approximately $10^{57}$ atoms, each with its own electron field envelope satisfying field closure $\det(e) \neq 0$. If each atom's field propagates at speed $c$ (locally enforcing closure), and the star's macroscopic field also propagates at speed $c$, how do these nested field envelopes remain consistent? Fields cannot coordinate instantaneously across the star's interior, yet field closure must hold continuously everywhere from the atomic scale ($\sim 10^{-10}$ m) to the stellar scale ($\sim 10^9$ m). The question is whether these are different equations at different scales, or the same equation universally:
\begin{equation}
\text{Atom: } \nabla^2\Phi_{\text{atom}} \stackrel{?}{=} \nabla^2\Phi_{\text{star}} \text{: Star} \stackrel{?}{=} \nabla^2\Phi_{\text{galaxy}} \text{: Galaxy}
\label{eq:scaling-question}
\end{equation}

The resolution lies in the non-linearity of the Field Origin. In tetrad gravity, the field equations are non-linear because they involve products of tetrad components: the closure condition $\det(e_\mu^a) \neq 0$ is a non-linear algebraic constraint. In bimetric gravity, the interaction potential $V(\mathcal{S}) = e_0(\mathcal{S}) + 16\,e_2(\mathcal{S}) + e_4(\mathcal{S})$ (Equation~\ref{eq:V-determined}) contains non-linear terms through the elementary symmetric polynomials $e_n$ of the mixing matrix $\mathcal{S} = \sqrt{g^{-1}f}$. This non-linearity is not a complication; it is the mechanism that allows field closure to be satisfied at every scale independently. Each gravitating system (atom, planet, star, galaxy) establishes its own field envelope with $k=4$ tetrahedral structure at that scale, without requiring global coordination with field envelopes at other scales. The topological constraint $k = d+1 = 4$ applies locally at each point in the continuous field, ensuring that closure holds everywhere regardless of the nesting hierarchy.

% ============================================================
% WHAT: MATHEMATICS (Topological Fixed Point and Conformal Symmetry)
% ============================================================

The key to understanding this scale independence is that $k=4$ is not a dynamical parameter that changes with energy or length scale; it is a \emph{topological constraint}. The simplex theorem requires $k = d+1$ tetrahedral faces in $d=3$ spatial dimensions. This constraint is geometric: it does not depend on the field strength, the mass density, or the characteristic scale of the system. In the language of renormalization group (RG) theory, $k=4$ is a ``fixed point'': as we examine the field at different scales $\lambda$, the value of $k$ remains constant. We quantify this with the beta function:
\begin{equation}
\beta(k) = \frac{\partial k}{\partial \ln \lambda} = 0
\label{eq:beta-function}
\end{equation}
where $\lambda$ is the length scale at which we probe the field. This beta function measures how a parameter changes as we zoom in or out; $\beta(k) = 0$ means $k$ does not ``run'' with scale. This is a \emph{topological fixed point}: the value $k=4$ is locked by the dimensionality of space, not by field dynamics. Unlike dynamical fixed points in quantum field theory (which arise from balance of competing interactions), this is a geometric constraint: the simplex theorem in three-dimensional space requires $k = d+1 = 4$ regardless of field strength or coupling constants.

At RG fixed points, field theories generically exhibit conformal symmetry: invariance under local rescalings of the metric. Specifically, under a Weyl (conformal) transformation:
\begin{equation}
g_{\mu\nu}(x) \to \tilde{g}_{\mu\nu}(x) = e^{2\omega(x)} g_{\mu\nu}(x)
\label{eq:weyl-transformation}
\end{equation}
where $\omega(x)$ is an arbitrary smooth function, distances are rescaled uniformly ($ds^2 \to e^{2\omega} ds^2$) but angles are preserved. If the field equations remain unchanged in form under such transformations, the theory is conformally symmetric, meaning it looks the same at every scale. This is exactly what we need to resolve the nested field envelope problem: if the Poisson equation is Weyl-invariant, then the same equation governs both an atom's field and a star's field, despite the vast difference in scale.

Now we can understand the factor $4\pi$ in the Poisson equation from this conformal perspective. The standard field theory derivation attributes $4\pi$ to the surface area of a sphere, $A = 4\pi r^2$. But surface area depends on radius $r$, which appears to introduce scale dependence. Our derivation is fundamentally different: the factor $4\pi$ comes from the product of discrete topology and continuous geometry:
\begin{equation}
\text{Field normalization} = k \times \pi = 4\pi
\label{eq:field-normalization}
\end{equation}

Field closure $\det(e) \neq 0$ requires $k=4$ tetrad basis vectors spanning the tangent space completely at every point; this is the tetrad formulation of GR's non-degeneracy condition~\cite{mtw1973,wald1984}. Complete spanning of four basis vectors defines a \emph{spherical simplex}: a regular tetrahedral ball at each vertex. The boundary of this simplex is a 2-sphere with total solid angle $4\pi$ steradians. This $4\pi$ is not derived from integration over radius $r$; it is intrinsic to the topology of spherical closure itself. The $k=4$ tetrahedral faces partition this measure: $4\pi = k \times \pi$, where each face carries $\pi$ steradians. Standard field theory arrives at $4\pi$ by integrating $\int\int r^2 \sin\theta\,d\theta\,d\phi$ over spherical shells; dual tetrad gravity recognizes that $4\pi$ is already encoded in the $k=4$ field topology that closure requires.

The critical distinction is:
\begin{align}
\text{Standard field theory:} \quad & 4\pi \text{ from } A = 4\pi r^2 \quad \text{(contains scale } r\text{)} \notag \\
\text{Dual tetrad gravity:} \quad & 4\pi = k \times \pi \text{ with } k=4 \quad \text{(manifestly scale-independent)}
\label{eq:4pi-distinction}
\end{align}
Because $k=4$ is a topological constant (the fixed point), the product $k\times\pi = 4\pi$ is automatically scale-invariant. There is no $r$ appearing anywhere in this expression; the factor $4\pi$ is a pure number determined by the topology of field closure, not by geometric integration over spherical shells. This normalization factor ensures scale invariance: the $k=4$ topology from tetrad closure~\cite{mtw1973,wald1984}, not the radius $r$ from geometric integration.

The same logic applies to the prefactor $17/13$. At the Field Origin ($s=0$, where both tetrahedra meet), the coupling polynomial evaluates to $f(k)|_{s=0} = k^2 + 1 = 17$. This integer counts the number of bimetric coupling modes: $k^2 = 16$ face-to-face channels between the two tetrahedra, plus the existence condition $\det(e) \neq 0$ (one additional constraint). In the spatial projection ($s=-1$, where we observe the field in $(3+1)$ spacetime), the polynomial evaluates to $f(k)|_{s=-1} = kd+1 = 13$ using $k^2 - k = kd$ for $k = d+1$. This counts the observable modes: $kd = 12$ spatial channels (three dimensions times four directions), plus the existence condition in projected space. The ratio of these two evaluations gives:
\begin{equation}
\frac{\text{Coupling capacity}}{\text{Observable capacity}}
= \frac{f(k)|_{s=0}}{f(k)|_{s=-1}}
= \frac{k^2+1}{kd+1}
= \frac{17}{13}
= 1.307692\ldots
\label{eq:coupling-ratio}
\end{equation}
Both the numerator and denominator are functions of the topological integers $k=4$ and $d=3$. Because these integers do not change with scale (they are geometric invariants), the ratio $17/13$ is also scale-invariant. It is not a phenomenological parameter fitted to match observations; it is a topological consequence of evaluating the polynomial at two structural states within the $k=4$ field architecture.

Combining the field normalization (\ref{eq:field-normalization}) with the coupling ratio (\ref{eq:coupling-ratio}) yields the complete Poisson equation that applies at all scales:
\begin{equation}
\nabla^2\Phi = \left[k\times\pi\right] \times \left[\frac{k^2+1}{kd+1}\right] \times G_{\text{bare}} \times \rho
\label{eq:poisson-general}
\end{equation}
Substituting $k=4$ and $d=3$:
\begin{equation}
\nabla^2\Phi = 4\pi \times \frac{17}{13} \times G_{\text{bare}} \times \rho
\label{eq:poisson-numeric}
\end{equation}
with numerical prefactor $4\pi \times (17/13) = 16.493362\ldots$. This equation is identical whether we are calculating the gravitational field of an electron ($\lambda \sim 10^{-13}$ m at the floor resolution), a planet ($\lambda \sim 10^6$ m), or a galaxy ($\lambda \sim 10^{21}$ m). All three factors (the normalization $4\pi$, the coupling ratio $17/13$, and the underlying $k=4$ tetrahedral structure) are determined by topological integers that do not change with scale. The only quantities that vary are the density $\rho$ and the characteristic scale $\lambda$ at which we evaluate the field, but these enter through the source term and boundary conditions, not through the differential operator or coupling constants.

For field object configurations (quantum mass body, e.g., an asteroid containing $10^{33}$ atoms), the non-linear Field Origin (established in both tetrad gravity through the non-linear closure constraint $\det(e) \neq 0$ and bimetric gravity through the non-linear interaction potential $V(\mathcal{S})$ containing elementary symmetric polynomials) ensures that each scale establishes $k=4$ closure independently without requiring global coordination. The field is not a simple average over constituent mass (this would destroy the topological structure). Rather, the field couples isomorphically to the object's configuration through the energy-momentum tensor $T_{\mu\nu}$, with the envelope-level field exhibiting $k=4$ topology when interacting with external fields. The physical object may be elongated and irregular, yet the envelope field (what external observers measure as the gravitational potential) satisfies the $k=4$ topological constraint because field closure $\det(e) \neq 0$ must hold everywhere. The result: $k=4$ at the envelope level ensures the Poisson equation (\ref{eq:poisson-numeric}) applies universally across all scales.

% ============================================================
% SO WHAT: IMPLICATIONS (Self-Similarity and Topological Origin)
% ============================================================

This is what ``self-similarity'' means in the context of gravitational fields: not that an atom's field is physically identical to a galaxy's field (they are obviously different), but that the \emph{equations} governing the fields are identical. The field topology ($k=4$ tetrahedral structure enforced by field closure $\det(e) \neq 0$) is the same at every scale because field closure is a local constraint, satisfied pointwise at each location in spacetime. An atom establishes field closure within its electron horizon; a planet establishes field closure within its gravitational radius; a galaxy establishes field closure within its virial radius. These are nested hierarchies, but at each level the constraint is the same: four linearly independent basis vectors (tetrads) must exist at every point. The simplex theorem guarantees this requires $k=4$ faces, and therefore the coupling polynomial $f(k) = 1 + k + k^2$ applies universally. Scale invariance follows directly from the topological origin of $k$.

The standard Poisson equation $\nabla^2\Phi = 4\pi G\rho$ appears to contain the same factor $4\pi$. Standard field theory derives this from the surface area of a sphere: integrate the gravitational field $\mathbf{g} = -\nabla\Phi$ over a closed surface surrounding a point mass, apply Gauss's law, and the surface area $A = 4\pi r^2$ cancels the $1/r^2$ decay of the field strength to yield $4\pi G M$ for the enclosed mass $M$. This derivation is mathematically correct. But it leaves a conceptual puzzle: the factor $4\pi$ appears to come from a geometric quantity (surface area) that depends on radius $r$, yet the final equation is universal across all scales. How does a scale-dependent quantity ($4\pi r^2$) produce a scale-independent result ($4\pi G\rho$)? The standard answer is that the $r^2$ in the surface area cancels the $1/r^2$ in the field, but this seems like a fortunate accident rather than a fundamental explanation.

Our topological derivation resolves this puzzle. The factor $4\pi = k\times\pi$ is manifestly scale-independent from the outset because $k=4$ is a topological integer (required by the simplex theorem and enforced by tetrad gravity's closure condition~\cite{mtw1973,wald1984}) that does not vary with $r$, $\lambda$, or any other scale parameter. Standard field theory arrives at the correct numerical value ($4\pi$) through geometric reasoning, but the deeper origin is topological: the four tetrahedral faces required by the simplex theorem in three-dimensional space, multiplied by the spherical geometry ($\pi$) of field propagation. This topological perspective explains \emph{why} the Poisson equation is universal: it is not because of a cancellation between scale-dependent geometric factors, but because the underlying field structure is determined by scale-independent topological constraints. The conformal symmetry at the $k=4$ fixed point ensures that Weyl transformations leave the equation covariant, which is the formal statement of scale invariance.

Section~\ref{subsec:prefactor} derived the integers $17$ and $13$ algebraically by evaluating the coupling polynomial $f(k) = 1 + sk + k^2$ at the Field Origin ($s=0$) and the spatial projection ($s=-1$). This section demonstrates that these integers, combined with the field normalization $k\times\pi$, produce the entire Poisson equation. Every component (the normalization $4\pi$, the prefactor $17/13$, and the effective gravitational constant $G_{\text{eff}} = (17/13)G_{\text{bare}}$) traces back to evaluating the polynomial within the constraint that $k = d+1 = 4$ is topologically fixed. This is the meaning of ``the field origin is the whole equation'': the Poisson equation is not constructed from disparate pieces (a geometric $4\pi$ from surface area, a dynamical $17/13$ from field coupling, and a separate constant $G$), but rather emerges as a unified structure from the polynomial evaluated at the Field Origin, with all factors determined by the same topological integers $k$ and $d$.

The scale-invariant form (\ref{eq:poisson-general}) can be written compactly:

\vspace{1.5em}
\noindent
\colorbox{resultgreen}{\parbox{\dimexpr\textwidth-2\fboxsep\relax}{%
\vspace{0.8em}
\centering
\textbf{\large The Poisson Equation from Field Closure}

\vspace{0.5em}
\boldmath
\begin{equation}
\nabla^2\Phi = k\pi \cdot \frac{k^2+1}{kd+1} \cdot G_{\text{bare}} \cdot \rho
\quad \text{with } k=4, \, d=3 \text{ (topological constants)}
\label{eq:poisson-boxed}
\end{equation}
\unboldmath

\vspace{0.3em}
\textbf{Scale-invariant from quantum to cosmic scales. Zero free parameters.}
\vspace{0.5em}}}
\vspace{1.5em}

The gravitational Poisson equation follows from field closure $\det(e) \neq 0$, the simplex theorem $k = d+1$, and the coupling polynomial $f(k) = 1 + k + k^2$ evaluated at structural states $s=0$ (Field Origin) and $s=-1$ (spatial projection). No parameters are fitted, adjusted, or borrowed from phenomenology. The conformal structure (topological fixed point, Weyl invariance, scale-independent factors) explains why this equation governs gravitational fields from the quantum regime (electron horizons) to the cosmic regime (galactic clusters), all without modification.
